% 中文学位论文/长篇报告骨架：XeLaTeX + ctexbook，字体集 fandol（离线字体包内置）。
% 结构：封面 → 中文摘要 → 英文摘要 → 目录 → 正文五章（绪论/相关工作/方法/实验/结论）
%       → 参考文献（BibTeX，GB/T 7714—2015）→ 附录 → 致谢。
% 这是通用骨架，不是任何学校的官方模板：正式提交前按学校/院系的格式要求调整
% 封面、页边距、字号与前置部分。章节较多时可把各章拆到 chapters/ 下的文件，
% 用 \include 命令引入。
\documentclass[UTF8,fontset=fandol,a4paper,zihao=-4,oneside,openany]{ctexbook}

\usepackage[top=2.5cm,bottom=2.5cm,left=3cm,right=2.5cm]{geometry}
\usepackage{amsmath,amssymb,amsthm}
\usepackage{booktabs}
\usepackage{graphicx}
% 参考文献采用 GB/T 7714—2015 顺序编码制（gbt7714 宏包 + gbt7714-numerical 样式），
% 这是国内学位论文通行的著录标准；默认上标引用，需要方括号正文引用时启用下一行。
\usepackage{gbt7714}
% \citestyle{numbers}
\usepackage{hyperref}
\hypersetup{colorlinks=true, linkcolor=black, citecolor=blue, urlcolor=blue}

% 参考文献标题不编号，但列入目录。
\renewcommand{\bibsection}{\chapter*{\bibname}\phantomsection\addcontentsline{toc}{chapter}{\bibname}}

% 定理类环境（中文名称，按章编号）；证明环境 proof 由 amsthm 提供，标题为"证明"。
\theoremstyle{plain}
\newtheorem{theorem}{定理}[chapter]
\newtheorem{lemma}[theorem]{引理}
\theoremstyle{definition}
\newtheorem{definition}[theorem]{定义}

\begin{document}

% ---------------- 封面（按学校要求替换）----------------
% 封面日期请写明答辩或提交时间；\today 会显示该版本内容首次构建的日期。
\pagenumbering{Alph}
\begin{titlepage}
  \centering
  \vspace*{3cm}
  {\zihao{-1}\heiti 论文题目\par}
  \vspace{1cm}
  {\zihao{3} 英文题目（English Title）\par}
  \vspace{3cm}
  {\zihao{4}
  \begin{tabular}{r@{：}l}
    作者姓名 & \underline{\makebox[6cm]{}} \\[1ex]
    指导教师 & \underline{\makebox[6cm]{}} \\[1ex]
    学科专业 & \underline{\makebox[6cm]{}} \\[1ex]
    所在学院 & \underline{\makebox[6cm]{}} \\[1ex]
    完成日期 & \underline{\makebox[6cm]{}} \\
  \end{tabular}\par}
  \vfill
\end{titlepage}

% ---------------- 前置部分（罗马数字页码）----------------
\frontmatter

\chapter*{摘\quad 要}
\addcontentsline{toc}{chapter}{摘要}
% 概述研究背景、问题、方法、主要结果与结论，篇幅按学校要求（常见为一页以内）。
在此填写中文摘要。

\bigskip
\noindent\textbf{关键词：}关键词一；关键词二；关键词三

\chapter*{Abstract}
\addcontentsline{toc}{chapter}{Abstract}
% English abstract; keep it consistent with the Chinese abstract.
Write the English abstract here.

\bigskip
\noindent\textbf{Keywords:} keyword one; keyword two; keyword three

\tableofcontents

% ---------------- 正文（阿拉伯数字页码）----------------
\mainmatter

\chapter{绪论}
% 研究背景与意义、国内外研究现状概述、研究内容与主要贡献、论文结构安排。
\section{研究背景与意义}
\section{研究内容与贡献}
\section{论文结构安排}

\chapter{相关工作}
% 按主题梳理已有研究并引用文献，指出尚未解决的问题。
\section{相关研究一}
\section{相关研究二}

\chapter{方法}
% 问题的形式化描述、核心定义与主要结论。以下环境仅示范写法，请替换为实际内容。
\section{问题定义}
\begin{definition}\label{def:main}
在此给出定义。
\end{definition}

编号公式示例（替换为实际公式）：
\begin{equation}\label{eq:main}
  f(x) = \sum_{i=1}^{n} w_i\, g_i(x) .
\end{equation}

\section{主要结论}
\begin{theorem}\label{thm:main}
在此陈述定理。
\end{theorem}

\begin{proof}
在此给出证明，可引用定义~\ref{def:main} 与式~\eqref{eq:main}。
\end{proof}

\chapter{实验}
% 数据、实验设置、评价指标、结果与分析。表中数据必须来自真实实验记录。
\section{实验设置}
\section{结果与分析}
\begin{table}[htbp]
  \centering
  \caption{实验结果（表头为示例，数据待填）}
  \label{tab:results}
  \begin{tabular}{lcc}
    \toprule
    方法     & 指标一 & 指标二 \\
    \midrule
    基线方法 & --     & --     \\
    本文方法 & --     & --     \\
    \bottomrule
  \end{tabular}
\end{table}

\chapter{结论}
% 总结全文工作与创新点，说明不足与后续研究方向。
\section{工作总结}
\section{不足与展望}

\bibliographystyle{gbt7714-numerical}
\bibliography{references}

% ---------------- 附录 ----------------
\appendix
\chapter{补充材料}
% 放置证明细节、补充实验、数据说明等不宜放入正文的内容。

% ---------------- 后置部分 ----------------
\backmatter
\chapter{致谢}
% 致谢内容。

\end{document}
